\documentclass[10pt,a4paper]{amsart}
\usepackage[margin=19mm]{geometry}
\usepackage{amsmath,amssymb,mathtools}
\usepackage[scr=rsfs,cal=euler]{mathalfa}
\usepackage{xfrac}
\usepackage[unicode,hidelinks,bookmarks=false]{hyperref}
\newcommand{\ie}{i.e.\ }
\newcommand{\cf}{cf.\ }
\newcommand{\resp}{respectively\ }
\newcommand{\Subsection}{\subsection}
\providecommand{\nobreakdash}{\nobreak\hbox{-}}
\setlength{\emergencystretch}{2em}
\multlinegap0pt
\newtheorem{thm}{Theorem}
\title{On Combinatorial Properties of the degenerate Krawtchouk Appell polynomials}
\author{Mohamed Abdelkader, Mohamed Rhaima}
\date{Source: arXiv:2605.16287v1 (CC BY 4.0)}
\begin{document}
\maketitle

\noindent\emph{Source and licence.} On Combinatorial Properties of the degenerate Krawtchouk Appell polynomials. Version arXiv:2605.16287v1 (CC BY 4.0), distributed by arXiv under the Creative Commons Attribution 4.0 International licence, http://creativecommons.org/licenses/by/4.0/, as recorded in the arXiv licence metadata for this exact version and shown on the abstract page https://arxiv.org/abs/2605.16287v1. Full licence notice: \texttt{licenses/arxiv-2605.16287-CC-BY-4.0.txt}.\par\smallskip

\noindent\textit{Editorial scope.} Counted unit: the article's thm def-lambda-mn (source0.tex lines 568--609). The article's complete original proof of it is delivered with it (source0.tex lines 610--751, one \texttt{proof} environment of 142 lines). No mathematics is written, repaired or re-derived: the statement and the proof are the article's own lines, in the article's own notation, and no retained line is substituted or re-flowed.  Retained uncounted as the source-local context the proof needs: lines 199--206; lines 299--304; lines 389--394; lines 562--567.  Nothing was omitted from any retained span, so the package carries no omission record.  No retained line contains a citation, so the wrapper carries no bibliography and no bibliography entry.  No retained line contains a figure environment, an image-inclusion command or drawing code, so no image is delivered and the package declares no asset dependency.  Editorial formatting only, in addition: an AMS \texttt{amsart} preamble that restates the article's own declarations and macros used by the retained lines, namely the environments \texttt{thm} on separate counters, together with the standard packages those lines need.  The retained environments print in this document as Theorem 1 in order of appearance, whereas the article prints the same items with the numbering of its own full document.  The complete native source of the article as distributed in the arXiv e-print for this exact version is bundled unchanged as \texttt{sources/200713/source0.tex}.
\begin{equation}\label{Degenerate-falling}
(\beta)_{k,\lambda}:=
\begin{cases}
1, & \text{if } k=0, \\
\beta(\beta- \lambda)(\beta- 2\lambda)\cdots (\beta-(k-1)\lambda), &
\text{if}\, k \in \mathbb{N}^{\ast},
\end{cases}
\end{equation}

\begin{equation}\label{m-th-Moments}
=\displaystyle{(1+\lambda r\log p)^{\frac{\beta}{\lambda}}
\sum^{n}_{k=1}S(m, k)\Big(\frac{q}{1+\lambda r\log
p}\Big)^{k}(\beta)_{k, \lambda}e^{\beta-\lambda
k}_{\lambda}\Big(\frac{q}{1+\lambda r\log p}\Big)},
\end{equation}

\begin{equation}\label{wick-exp}
\mathcal{N}\!\exp^{\beta}_{\lambda}(z,
x):=\displaystyle{\frac{e^{xz}}{\mathcal{L}_{\mathcal{N}\mathcal{B}^{(\lambda)}_{p,
r}}(z)}
=\frac{e^{xz}}{e^{\beta}_{\lambda}\big[r\log(\frac{p}{1-qe^{z}})\big]}}.
\end{equation}

\begin{equation}\label{Exp-alpha-1}
\mathcal{N}\!\exp^{\beta}_{\lambda}(z,
x)=\displaystyle{\frac{e^{xz}}{\mathcal{L}_{\mathcal{N}\mathcal{B}^{(\lambda)}_{p,
r}}(z)}= \sum^{\infty}_{n=0}\frac{z^{n}}{n!}P^{(\beta)}_{n,
\lambda}(x)},\quad z\in \mathcal{U}.
\end{equation}

\begin{thm}
The polynomials $\mathcal{K}^{(\beta)}_{n, \lambda}$ satisfy the
following useful combinatorial properties:
\begin{itemize}
\item [\textbf{(P1)}]
$\mathcal{K}^{(\beta)}_{n, \lambda}(x)=\displaystyle{\sum^{n}_{m=0}
\frac{\varpi(m, n)}{m!}P^{(\beta)}_{m, \lambda}(x)}$, where
$\varpi(0, 0):=1$ and
\begin{equation}\label{def-lambda-mn}
\varpi(m,
n):=\displaystyle{\sum_{i_{1}+\ldots+i_{m}=n}(-1)^{n+m}n!\prod^{m}_{j=1}\Big(\frac{1-q^{i_{j}}}{i_{j}}\Big)}.
\end{equation}
\item [\textbf{(P2)}]
$x^{n}=\displaystyle{\sum^{n}_{k=0}\sum^{k}_{m=0}\binom{n}{k}
\frac{\varrho(m, k)}{m!}\mathcal{K}^{(\beta)}_{m,
\lambda}(x)M_{\lambda}(n-k)}$, where $M_{\lambda}(k)$ is the $k$-th
moment given by equation \eqref{m-th-Moments} and
$$\begin{cases}
\varrho(0, 0):=1\\
\varrho(m,
k):=\displaystyle{\sum_{i_{1}+\cdots+i_{m}=k}\frac{k!}{i_{1}!\ldots
i_{m}!}\prod^{m}_{j=1}\frac{d^{i_{j}}}{dz^{i_{j}}}\Big[\frac{e^{z}-1}{1-qe^{z}}\Big]_{z=0}
}.
\end{cases}
$$
\item [\textbf{(P3)}]
$\mathcal{K}^{(\beta)}_{n,
\lambda}(x+y)=\displaystyle{\sum_{k+l+m=n}
\frac{n!}{k!l!m!}\mathcal{K}^{(\beta)}_{k,
\lambda}(x)\mathcal{K}^{(\beta)}_{l,
\lambda}(x)(e^{\beta}_{\lambda}\circ \xi_{q})^{(m)}(0)}$.
\item [\textbf{(P4)}]
$\mathcal{K}^{(\beta)}_{n,
\lambda}(x+y)=\displaystyle{\sum^{n}_{k=0}\binom{n}{k}\mathcal{K}^{(\beta)}_{k,
\lambda}(x)[y]_{n-k}}$,
\end{itemize}
where
\begin{equation}\label{def-of-[y]m}
[y]_{n}:=\displaystyle{\sum^{n}_{k=0}\binom{n}{k}q^{k}(y)_{1,
n-k}(-y)_{1, k}}.
\end{equation}
\end{thm}

\begin{proof}
\textbf{(P1)}
Firstly by using equation (\ref{Exp-alpha-1}), we get
$$\Psi_{\lambda}(x, z)=\displaystyle{\frac{e^{x\vartheta(z)}}{\mathcal{L}_{\mathcal{NB}^{(\lambda)}_{p, r}}(\vartheta(z))}
=\sum^{\infty}_{m=0}\frac{\vartheta(z)^{m}}{m!}A^{(\beta)}_{m,
\lambda}},$$ for any $z\in \mathcal{U}$. It is obvious that
$$\vartheta(z)=\log\Big(\frac{1+z}{1+qz}\Big)=\displaystyle{\sum^{\infty}_{k=0}\frac{z^{k}}{k!}\eta_{k}},$$
with
$$
\eta_{k}=\displaystyle{(-1)^{k+1}(k-1)!(1-q^{k}) \quad\forall
k>0,\quad \eta_{0}=0}.
$$
Therefore with the help of Eq. \eqref{def-lambda-mn} we have\\

$\displaystyle{\sum^{\infty}_{n=0}\frac{z^{n}}{n!}\mathcal{K}^{(\beta)}_{n,
\lambda}(x)}$
$$\begin{array}{lll}
\qquad&=&P^{(\beta)}_{0,
\lambda}+\displaystyle{\sum^{\infty}_{m=1}\frac{1}{m!}\Big(\sum^{\infty}_{k=0}
\frac{z^{k}}{k!}\eta_{k}\Big)^{m}P^{(\beta)}_{m, \lambda}(x)}\\
&=&P^{(\beta)}_{0,
\lambda}+\displaystyle{\sum^{\infty}_{m=1}\frac{1}{m!}\sum^{\infty}_{n=m=0}\frac{z^{n}}{n!}
\sum_{i_{1}+\cdots+i_{m}=n}\frac{n!}{i_{1}!\cdots i_{m}!}\prod^{m}_{j=1}\eta_{i_{j}}P^{(\beta)}_{m, \lambda}(x)}\\
&=&P^{(\beta)}_{0,
\lambda}+\displaystyle{\sum^{\infty}_{m=1}\frac{1}{m!}\sum^{\infty}_{n=m}\frac{z^{n}}{n!}
\sum_{i_{1}+\cdots+i_{m}=n}(-1)^{n+m}n!\prod^{m}_{j=1}\Big(\frac{1-q^{i_{j}}}{i_{j}}\Big)P^{(\beta)}_{m, \lambda}(x)}\\
&=&\displaystyle{\sum^{\infty}_{n=1}\frac{z^{n}}{n!}\sum^{n}_{m=0}\frac{\varpi(m,
n)}{m!}P^{(\beta)}_{m, \lambda}(x)}.
\end{array}$$
Hence we obtain
$$\displaystyle{\mathcal{K}^{(\beta)}_{n, \lambda}(x)=\displaystyle{\sum^{n}_{m=0}
\frac{\varpi(m, n)}{m!}P^{(\beta)}_{m, \lambda}(x)}}.$$
\begin{itemize}
\item [\textbf{(P2)}]
Let $\zeta$ be the inverse function of $\vartheta$. Then one can see
that
$$\zeta(z)=\displaystyle{\frac{e^{z}-1}{1-qe^{z}}=\sum^{\infty}_{k=0}\kappa_{k}\frac{z^{k}}{k!},}$$
where
$$\kappa_{k}=\displaystyle{\frac{d^{k}}{dz^{k}}\Big[\frac{e^{z}-1}{1-qe^{z}}\Big]_{z=0}}.$$
By using equation (\ref{wick-exp}), we get
$$\displaystyle{\frac{e^{xz}}{\mathcal{L}_{\mathcal{NB}^{(\lambda)}_{p, r}}(z)}=
\frac{e^{xz}}{e^{\beta}_{\lambda}\Big(r\log(\frac{p}{1-qe^{z}})\Big)}=
\sum^{\infty}_{m=0}\frac{\zeta(z)^{m}}{m!}\mathcal{K}^{(\beta)}_{m,
\lambda}(x)},\quad z\in \mathcal{U}.$$
So we have\\

$\displaystyle{\sum^{\infty}_{n=0}\frac{z^{n}}{n!}A^{(\beta)}_{n,
\lambda}(x)}=\mathcal{K}^{(\beta)}_{0, \lambda}(x)
+\displaystyle{\sum^{\infty}_{m=1}\frac{1}{m!}\Big(\sum^{\infty}_{k=1}\kappa_{k}\frac{z^{k}}{k!}\Big)^{m}\mathcal{K}^{(\beta)}_{m,
\lambda}(x)} $
$$\begin{array}{lll}
&=&\quad \mathcal{K}^{(\beta)}_{0,
\lambda}(x)+\displaystyle{\sum^{\infty}_{m=1}
\frac{1}{m!}\sum^{\infty}_{n=m}\frac{z^{n}}{n!}\sum_{i_{1}+\ldots+i_{m}=n}\frac{n!}{i_{1}!\ldots
i_{m}!}\prod^{m}_{j=1}\kappa_{i_{j}}
\mathcal{K}^{(\beta)}_{m, \lambda}(x)}\\
&=&\displaystyle{\sum^{\infty}_{n=0}\frac{z^{n}}{n!}
\sum^{n}_{m=0}\frac{\varrho(m, n)}{m!}\mathcal{K}^{(\beta)}_{m,
\lambda}(x)}.
\end{array}$$
Thus, we have
\begin{equation}\label{Aalphan}
P^{(\beta)}_{n,
\lambda}(x)=\displaystyle{\sum^{n}_{m=0}\frac{\varrho(m,
n)}{m!}\mathcal{K}^{(\beta)}_{m, \lambda}(x)}.
\end{equation}
On the other hand since
$$e^{xz}=\mathcal{N}\!\exp^{\beta}_{\lambda}(z, x)\mathcal{L}_{\mathcal{NB}^{(\lambda)}_{p, r}}(z),$$
we get
$$\begin{array}{lll}
\displaystyle{\sum^{\infty}_{n=0}\frac{x^{n}}{n!}z^{n}}&=&\displaystyle{\sum^{\infty}_{k=0}\frac{z^{k}}{k!}P^{(\beta)}_{k,
\lambda}(x)
\sum^{\infty}_{k=0}\frac{z^{k}}{k!}M_{\lambda}(k)}\\
&=&\displaystyle{\sum^{\infty}_{n=0}\frac{z^{n}}{n!}\Big(\sum^{n}_{k=0}\binom{n}{k}
P^{(\beta)}_{k, \lambda}(x)M_{\lambda}(n-k)\Big)}.
\end{array}$$
Hence with the help of Eq. \eqref{Aalphan}, we obtain
\begin{eqnarray}\label{eq-power-xn}
x^{n}&=&\displaystyle{\sum^{n}_{k=0}\binom{n}{k}
P^{(\beta)}_{k, \lambda}(x)M_{\lambda}(n-k)}\\
&=&\displaystyle{\sum^{n}_{k=0}\sum^{k}_{m=0}
\binom{n}{k}\frac{\varrho(m, k)}{m!}\mathcal{K}^{(\beta)}_{m,
\lambda}(x)M_{\lambda}(n-k)}.\nonumber
\end{eqnarray}
\item [\textbf{(P3)}]
By using the fact that
$$\mathcal{N}\!\exp^{\beta}_{\lambda}(\vartheta(z), x+y)=\mathcal{N}\!\exp^{\beta}_{\lambda}(\vartheta(z), x)
\mathcal{N}\!\exp^{\beta}_{\lambda}(\vartheta(z),
y)\mathcal{L}_{\mathcal{NB}^{(\lambda)}_{p, r}}(\vartheta(z)),$$ and
$$
\mathcal{L}_{\mathcal{NB}^{(\lambda)}_{p,
r}}(\vartheta(z))=e^{\beta}_{\lambda}\big(r\log(1+qz)\big)
=\displaystyle{\sum^{\infty}_{k=0}\frac{(e^{\beta}_{\lambda}\circ
\xi_{q})^{(k)}(0)}{k!}z^{k}},$$ we obtain
$$\begin{array}{lll}
\displaystyle{\sum^{\infty}_{n=0}\frac{z^{n}}{n!}\mathcal{K}^{(\beta)}_{n,
\lambda}(x+y)}
&=&\displaystyle{\sum^{\infty}_{k=0}\frac{z^{k}}{k!}\mathcal{K}^{(\beta)}_{k,
\lambda}(x)
\sum^{\infty}_{l=0}\frac{z^{l}}{l!}\mathcal{K}^{(\beta)}_{l,
\lambda}(y)
\sum^{\infty}_{m=0}\frac{z^{m}}{m!}\mathcal{K}^{(\beta)}_{m, \lambda}(x)}\\
&=&\displaystyle{\sum^{\infty}_{n=0}\frac{z^{n}}{n!}\sum_{k+l+m=n}
\frac{n!}{k!l!m!}\mathcal{K}^{(\beta)}_{k,
\lambda}(x)\mathcal{K}^{(\beta)}_{l,
\lambda}(x)(e^{\beta}_{\lambda}\circ \xi_{q})^{(m)}(0)}.
\end{array}$$
This gives the statement.
\item [\textbf{(P4)}]
Recall that, for $n\in \mathbb{N}$ and  $y\in \mathbb{C}$, the
generating function of the classical falling factorials
$(y)_{n}:=(y)_{1, n}$ obtained by \eqref{Degenerate-falling} with
$\lambda=1$ is
$$
\displaystyle{\sum^{\infty}_{n=0}\frac{z^{n}}{n!}(y)_{n}=\exp[y\log(1+z)]}.
$$
Then one can directly see that
\begin{eqnarray}\label{Taylor-exp-exp-v}
\exp(y\vartheta(z))&=&\exp(y\log(1+z)-y\log(1+qz))\nonumber\\
&=&\displaystyle{\sum^{\infty}_{n=0}[y]_{n}\frac{z^{n}}{n!}},
\end{eqnarray}
where
$$
[y]_{n}:=\displaystyle{\sum^{n}_{k=0}\binom{n}{k}q^{k}(y)_{n-k}(-y)_{k}}.
$$
It is clear that
$$\mathcal{N}\!\exp^{\beta}_{\lambda}(\vartheta(z), x+y)=\mathcal{N}\!\exp^{\beta}_{\lambda}(\vartheta(z), x)\exp(y\vartheta(z)).$$
Then with the help of equation \eqref{Taylor-exp-exp-v} we get
$$\begin{array}{lll}
\displaystyle{\sum^{\infty}_{n=0}\frac{z^{n}}{n!}\mathcal{K}^{(\beta)}_{n,
\lambda}(x+y)}&=&
\displaystyle{\sum^{\infty}_{k=0}\frac{z^{k}}{k!}\mathcal{K}^{(\beta)}_{k, \lambda}(x)\sum^{\infty}_{k=0}\frac{z^{k}}{k!}[y]_{k}}\\
&=&\displaystyle{\sum^{\infty}_{n=0}\frac{z^{n}}{n!}\sum_{k+m=n}\frac{n!}{k!m!}\mathcal{K}^{(\beta)}_{k, \lambda}(x)[y]_{m}}\\
&=&\displaystyle{\sum^{\infty}_{n=0}\frac{z^{n}}{n!}\sum^{n}_{k=0}
\binom{n}{k}\mathcal{K}^{(\beta)}_{k, \lambda}(x)[y]_{n-k}}.
\end{array}$$
Hence we deduce the expansion of $\mathcal{K}^{(\beta)}_{n,
\lambda}(x+y)$ given by
$$\mathcal{K}^{(\beta)}_{n, \lambda}(x+y)=\displaystyle{\sum^{n}_{k=0}
\binom{n}{k} \mathcal{K}^{(\beta)}_{k, \lambda}(x)[y]_{n-k}}.$$
\end{itemize}
\end{proof}

\end{document}
